\textbf{Phần 1. Trắc nghiệm lựa chọn}
\setcounterpageref{ex}{0}

\begin{ex} %[5X11N66]
Khi thực hiện khảo sát sở thích của các bạn trong lớp, phương pháp thu thập dữ liệu phù hợp là
\choice
{Đo lường bằng thước kẹp}
{Cân bằng cân điện tử}
{\True Phát phiếu khảo sát hoặc phỏng vấn trực tiếp}
{Bấm giờ bằng đồng hồ}
\loigiai{
Khảo sát sở thích thường dùng phiếu hỏi hoặc phỏng vấn trực tiếp.\\
\textbf{Đáp án của câu hỏi là C}
}
\end{ex}

\begin{ex} %[5X11N66]
Dạng biểu đồ nào thích hợp nhất để so sánh trực quan số lượng học sinh thích các môn thể thao khác nhau?
\choice
{\True Biểu đồ cột hoặc biểu đồ hình quạt tròn}
{Biểu đồ đường gấp khúc}
{Sơ đồ hình cây}
{Biểu đồ phân tán}
\loigiai{
Biểu đồ cột hoặc biểu đồ hình quạt tròn rất trực quan để so sánh cơ cấu và số lượng.\\
\textbf{Đáp án của câu hỏi là A}
}
\end{ex}

\begin{ex} %[5X11H66]
Khảo sát về loại trái cây yêu thích của $40$ bạn học sinh lớp 5B: Xoài ($14$ bạn), Cam ($10$ bạn), Ổi ($8$ bạn), Dưa hấu ($8$ bạn). Loại quả được nhiều bạn yêu thích nhất là
\choice
{Cam}
{\True Xoài}
{Ổi}
{Dưa hấu}
\loigiai{
Xoài có $14$ bạn, nhiều nhất.\\
\textbf{Đáp án của câu hỏi là B}
}
\end{ex}

\begin{ex} %[5X11H66]
Bảng khảo sát số cuốn sách đọc được trong tuần của một nhóm bạn:
- $1$ cuốn: $3$ bạn
- $2$ cuốn: $5$ bạn
- $3$ cuốn: $2$ bạn
Tổng số cuốn sách nhóm bạn đã đọc trong tuần là
\choice
{$10$ cuốn}
{$20$ cuốn}
{\True $19$ cuốn}
{$15$ cuốn}
\loigiai{
Tổng số sách: $1 \times 3 + 2 \times 5 + 3 \times 2 = 3 + 10 + 6 = 19$ (cuốn).\\
\textbf{Đáp án của câu hỏi là C}
}
\end{ex}

\begin{ex} %[5X11H66]
Khảo sát thời gian xem ti vi mỗi ngày của các bạn trong tổ: $30$ phút, $45$ phút, $60$ phút, $45$ phút, $30$ phút, $90$ phút. Thời gian xem ti vi trung bình của một bạn trong tổ là
\choice
{$45$ phút}
{$55$ phút}
{$60$ phút}
{\True $50$ phút}
\loigiai{
Tổng thời gian: $30 + 45 + 60 + 45 + 30 + 90 = 300$ phút.\\
Trung bình: $300 : 6 = 50$ phút.\\
\textbf{Đáp án của câu hỏi là D}
}
\end{ex}

\begin{ex} %[5X11H66]
Khi biểu diễn số liệu học sinh tham gia các hoạt động ngoại khóa: Văn nghệ chiếm $40\%$, Thể thao chiếm $60\%$ lên biểu đồ hình quạt tròn, phần thể thao chiếm
\choice
{Ít hơn một nửa hình tròn}
{\True Nhiều hơn một nửa hình tròn}
{Đúng một nửa hình tròn}
{Toàn bộ hình tròn}
\loigiai{
$60\% > 50\%$ nên phần thể thao chiếm nhiều hơn một nửa hình tròn.\\
\textbf{Đáp án của câu hỏi là B}
}
\end{ex}

\begin{ex} %[5X11V66]
Kết quả đo chiều cao của $5$ học sinh: $138\text{ cm}, 142\text{ cm}, 140\text{ cm}, 136\text{ cm}, 144\text{ cm}$. Chiều cao trung bình của $5$ học sinh đó là
\choice
{\True $140\text{ cm}$}
{$139\text{ cm}$}
{$141\text{ cm}$}
{$142\text{ cm}$}
\loigiai{
Tổng chiều cao: $138 + 142 + 140 + 136 + 144 = 700\text{ cm}$.\\
Trung bình: $700 : 5 = 140\text{ cm}$.\\
\textbf{Đáp án của câu hỏi là A}
}
\end{ex}

\begin{ex} %[5X11V66]
Khảo sát số thành viên trong gia đình của $30$ bạn học sinh thu được:
- Gia đình có $3$ người: $6$ bạn
- Gia đình có $4$ người: $18$ bạn
- Gia đình có $5$ người: $6$ bạn
Tỉ số phần trăm số bạn có gia đình gồm $4$ người là
\choice
{$50\%$}
{$40\%$}
{$70\%$}
{\True $60\%$}
\loigiai{
$18 : 30 = 0{,}6 = 60\%$.\\
\textbf{Đáp án của câu hỏi là D}
}
\end{ex}

\begin{ex} %[5X11V66]
Một nhóm học sinh đo nhiệt độ lúc $12$ giờ trưa trong $5$ ngày liên tiếp: $28^\circ\text{C}, 30^\circ\text{C}, 32^\circ\text{C}, 31^\circ\text{C}, 29^\circ\text{C}$. Nhiệt độ trung bình trong $5$ ngày là
\choice
{$29^\circ\text{C}$}
{\True $30^\circ\text{C}$}
{$31^\circ\text{C}$}
{$30{,}5^\circ\text{C}$}
\loigiai{
Tổng nhiệt độ: $28 + 30 + 32 + 31 + 29 = 150^\circ\text{C}$.\\
Trung bình: $150 : 5 = 30^\circ\text{C}$.\\
\textbf{Đáp án của câu hỏi là B}
}
\end{ex}

\begin{ex} %[5X11V66]
Trong dự án “Nuôi heo đất giúp bạn vượt khó”, bốn tổ của lớp 5A quyên góp được số tiền lần lượt là: $250\,000$ đồng, $300\,000$ đồng, $280\,000$ đồng, $370\,000$ đồng. Trung bình mỗi tổ quyên góp được số tiền là
\choice
{\True $300\,000$ đồng}
{$280\,000$ đồng}
{$320\,000$ đồng}
{$350\,000$ đồng}
\loigiai{
Tổng tiền: $250 + 300 + 280 + 370 = 1\,200$ nghìn đồng.\\
Trung bình: $1\,200 : 4 = 300$ nghìn đồng $= 300\,000$ đồng.\\
\textbf{Đáp án của câu hỏi là A}
}
\end{ex}

\begin{ex} %[5X11V66]
Thống kê thời gian tự học ở nhà (phút/ngày) của $10$ học sinh: $45, 60, 50, 60, 45, 90, 60, 75, 60, 45$. Tỉ số phần trăm học sinh dành đúng $60$ phút tự học là
\choice
{$30\%$}
{$50\%$}
{$20\%$}
{\True $40\%$}
\loigiai{
Có 4 bạn dành đúng 60 phút $\Rightarrow 4 : 10 = 40\%$.\\
\textbf{Đáp án của câu hỏi là D}
}
\end{ex}

\begin{ex} %[5X11C66]
Một cuộc khảo sát $100$ người tiêu dùng về $3$ nhãn hiệu sữa A, B, C cho thấy: $45$ người thích nhãn hiệu A, $35$ người thích nhãn hiệu B, số còn lại thích nhãn hiệu C. Khi vẽ biểu đồ quạt tròn, góc ở tâm biểu diễn phần số người thích nhãn hiệu C là bao nhiêu độ?
\choice
{$60^\circ$}
{$80^\circ$}
{\True $72^\circ$}
{$90^\circ$}
\loigiai{
Số người thích C: $100 - (45 + 35) = 20$ người ($20\%$).\\
Góc ở tâm: $360^\circ \times 20\% = 72^\circ$.\\
\textbf{Đáp án của câu hỏi là C}
}
\end{ex}

\textbf{Phần 2. Trắc nghiệm đúng sai}
\setcounterpageref{ex}{0}

\begin{ex} %[5X11H66]
Bảng khảo sát món tráng miệng ưa thích của $40$ học sinh: Kem ($16$ bạn), Chè ($12$ bạn), Trái cây ($8$ bạn), Sữa chua ($4$ bạn). Xét tính đúng sai:
\choiceTF
{\True Tỉ số phần trăm học sinh thích Kem là $40\%$.}
{\True Tỉ số phần trăm học sinh thích Chè là $30\%$.}
{Số học sinh thích Sữa chua chiếm $5\%$.}
{\True Số học sinh thích Sữa chua chiếm $10\%$.}
\loigiai{
\begin{itemize}
    \item Kem: $16 : 40 = 40\%$. Đúng.
    \item Chè: $12 : 40 = 30\%$. Đúng.
    \item Sữa chua: $4 : 40 = 10\% \neq 5\%$. Ý 3 Sai, Ý 4 Đúng.
\end{itemize}
\begin{itemize}
    \item a) \textbf{Đúng}.
    \item b) \textbf{Đúng}.
    \item c) \textbf{Sai}.
    \item d) \textbf{Đúng}.
\end{itemize}
}
\end{ex}

\begin{ex} %[5X11V66]
Một nhóm bạn đo lượng nước uống mỗi ngày của $5$ thành viên: $1{,}5\text{ l}, 1{,}8\text{ l}, 2{,}0\text{ l}, 1{,}7\text{ l}, 2{,}0\text{ l}$. Xét tính đúng sai:
\choiceTF
{\True Tổng lượng nước $5$ bạn uống trong ngày là $9{,}0\text{ l}$.}
{\True Lượng nước uống trung bình của mỗi bạn là $1{,}8\text{ l}$.}
{Có $1$ bạn uống lượng nước nhiều hơn mức trung bình.}
{\True Có $2$ bạn uống lượng nước nhiều hơn mức trung bình ($2{,}0\text{ l} > 1{,}8\text{ l}$).}
\loigiai{
\begin{itemize}
    \item Tổng: $1{,}5 + 1{,}8 + 2{,}0 + 1{,}7 + 2{,}0 = 9{,}0\text{ l}$. Đúng.
    \item Trung bình: $9{,}0 : 5 = 1{,}8\text{ l}$. Đúng.
    \item Có 2 bạn uống 2,0 l $> 1,8$ l. Ý 3 Sai, Ý 4 Đúng.
\end{itemize}
\begin{itemize}
    \item a) \textbf{Đúng}.
    \item b) \textbf{Đúng}.
    \item c) \textbf{Sai}.
    \item d) \textbf{Đúng}.
\end{itemize}
}
\end{ex}

\textbf{Phần 3. Trả lời ngắn}
\setcounterpageref{ex}{0}

\begin{ex} %[5X11H66]
Khảo sát số thành viên của $5$ gia đình: $3, 4, 5, 4, 4$. Tính số thành viên trung bình của mỗi gia đình.
\shortans{4}
\loigiai{
\textbf{Đáp số:} $4$\\
$(3 + 4 + 5 + 4 + 4) : 5 = 20 : 5 = 4$.
}
\end{ex}

\begin{ex} %[5X11H66]
Một bạn gieo đồng xu $25$ lần thấy mặt Sấp xuất hiện $13$ lần. Tính số lần xuất hiện mặt Ngửa.
\shortans{12}
\loigiai{
\textbf{Đáp số:} $12$\\
$25 - 13 = 12$.
}
\end{ex}

\begin{ex} %[5X11V66]
Một lớp có $35$ học sinh. Qua khảo sát có $28$ bạn thích xem hoạt hình. Tính tỉ số phần trăm học sinh thích xem hoạt hình. (Chỉ ghi số).
\shortans{80}
\loigiai{
\textbf{Đáp số:} $80$\\
$28 : 35 = 0{,}8 = 80\%$.
}
\end{ex}

\begin{ex} %[5X11V66]
Khảo sát số cây hoa trồng trong ban công của $4$ nhà: $8, 12, 15, 9$. Tính số cây hoa trung bình mỗi nhà trồng được.
\shortans{11}
\loigiai{
\textbf{Đáp số:} $11$\\
$(8 + 12 + 15 + 9) : 4 = 44 : 4 = 11$.
}
\end{ex}

\begin{ex} %[5X11V66]
Một nhóm học sinh quyên góp sách cũ: Tổ 1 góp $45$ cuốn, Tổ 2 góp $55$ cuốn, Tổ 3 góp $60$ cuốn, Tổ 4 góp $40$ cuốn. Tính số sách trung bình mỗi tổ quyên góp được.
\shortans{50}
\loigiai{
\textbf{Đáp số:} $50$\\
$(45 + 55 + 60 + 40) : 4 = 200 : 4 = 50$.
}
\end{ex}

\begin{ex} %[5X11C66]
Điểm kiểm tra của $4$ bạn là: $7, 8, 9, 8$. Hỏi bạn thứ năm cần đạt bao nhiêu điểm để điểm trung bình của cả nhóm $5$ bạn đạt $8{,}4$?
\shortans{10}
\loigiai{
\textbf{Đáp số:} $10$\\
Tổng điểm cần: $8{,}4 \times 5 = 42$.\\
Tổng 4 bạn đầu: $7 + 8 + 9 + 8 = 32$.\\
Điểm bạn thứ năm: $42 - 32 = 10$.
}
\end{ex}
